\subsection{Matched design pipelines on distinct geometry}
\label{sec:submissionstudy}

We next isolate three questions that a cavity-count comparison alone cannot
answer: whether objective-aware nucleation improves subsequent boundary
optimization; whether accessible-channel improvements survive fixed material
and power budgets; and whether the shrinking-target comparison transfers to
actual three-dimensional application meshes. The finned housing, branching
manifold, and ribbed annular spreader in \cref{fig:submissiongeometries} have
narrow gaps, branching paths, or curved transport and a thin central bridge.
Their protected sensors and complete compact heat-source supports remain
unchanged under every design. \Cref{app:submissionstudy} specifies the physics,
normalizations, geometry constraints, and numerical checks.



\paragraph{Nucleation followed by Differential WoS.}
All methods place six disjoint spheres of radius $0.065$ from the same paired
candidate pools, then optimize their centers using a common specialized
implementation of Differential WoS~\cite{miller2024differential}.
TWoS and deterministic topological derivatives use eight refinement rounds;
hottest-point and random initializations also receive rounds nine through
sixteen. Two complete deterministic baselines use longest-axis grid counts
49 and 65. A nearest-sensor control replaces the PDE score by Euclidean
distance to the protected sensor and receives eight refinement rounds. The curves in \cref{fig:submissionpipeline} report complete design
costs at independently evaluated checkpoints. This comparison isolates the
value of initialization inside a shared optimizer; it does not compare against
the original authors' full application or assert that Differential WoS cannot
change topology.

Across five paired seeds, TWoS followed by eight refinement rounds yields
mean sensor reductions of $65.0\%$, $16.4\%$, and $11.4\%$ on the housing,
branching manifold, and annular spreader. The primary multigrid selector followed
by the same refinement yields $65.1\%$, $16.3\%$, and $11.7\%$.
Hottest-point initialization with sixteen rounds yields $42.7\%$, $0.14\%$,
and $0.01\%$, respectively. Both TD selectors perform similarly in these
comparisons; their numerical costs and the additional controls are reported
in the full time--quality curves and \cref{tab:submissionpipeline}.
The nearest-sensor control yields $63.6\%$, $17.3\%$, and $12.2\%$
across the same three objects. On the housing it requires $7.08$ s, compared
with $11.77$ s for TWoS.
Thus the large contrast with hottest-point placement does not establish a
large benefit over every informed heuristic. The cutaways show this stronger
geometric control alongside TWoS and multigrid; the plots retain all methods.
The bands describe uncertainty in measured mean cooling, not equivalence tests
or guarantees for a new geometry.


\paragraph{Accessible channels at matched volume and fixed power.}
Six parallel open ball-union channels target total removed volume $0.0168$.
Their common radius is recalibrated after each insertion, and independent
scrambled-Sobol samples audit the final volume. All heat supports are protected,
so their exact input power is invariant (\cref{prop:protectedpower}).
\Cref{fig:submissionchannels} reports these finite designs. A pointwise
axis score initializes the channels; it is not a topological derivative of
an entire channel, and the radius calibration is numerical rather than certified.
At this matched budget, TWoS yields mean reductions of $37.8\%$, $2.29\%$,
and $1.75\%$ across the three objects, compared with $42.1\%$, $2.59\%$,
and $1.93\%$ for the primary multigrid selector. On the housing, nearest-sensor
placement yields $39.4\%$, and the cheaper multigrid grid attains the same
mean reduction as the primary grid. Thus this finite-channel heuristic does
not dominate the informed controls. \Cref{tab:submissionchannels} reports all
methods and uncertainty intervals; the independent volume and reference audits
remain separate from those sampling intervals.

\paragraph{Variance, bias, and query cost.}
\Cref{fig:submissionestimators} compares direct scores, common-random-number
finite-cavity subtraction, a stronger stopped-path harmonic-correction
baseline, and multigrid at four resolutions on eight fixed sites per object.
Sixteen independent repeats resolve sampling variability. Relative variance
about each method's mean is shown separately from error to a high-budget
infinitesimal reference: finite-cavity methods retain finite-radius bias,
while multigrid retains discretization error. The reference itself has
sampling and termination error. Bands around variance describe uncertainty
in the estimated variance, not a standard-deviation envelope of field values.
\Cref{app:submissionstudy} gives the stronger finite-correction identity and
its step-by-step proof; detailed numeric summaries accompany the source.
At radius $0.01$ and 2,048 paths per site, the stronger harmonic-correction
baseline has $35.7\times$ and $13.8\times$ the measured direct variance on
the housing and branching object. These are point estimates, not confidence
bounds. On the annular object, all recorded harmonic corrections vanish at
radii $0.01$ and $0.02$; variance is unresolved, and collapsed bootstrap
intervals are suppressed. At larger radii the stronger correction can have
competitive measured variance. Direct TWoS relative RMS error at this budget
is $14.9\%$, $25.2\%$, and $146\%$ across the three cases, while the
reference's relative standard-error norms are $1.6\%$, $4.8\%$, and
$16.2\%$. The annular sparse-query configuration is poorly resolved at these
budgets; it is not evidence of uniformly accurate pointwise estimation.

The maximum independent channel-volume audit deviation is $0.97\%$ of the
target. Finer housing checks change by at most $0.44$ percentage points from
grid 193 to 257, but MC--grid257 differences still reach $2.38$ points across
the checked first-seed states. Sampling intervals do not absorb this numerical
reference sensitivity (\cref{tab:submissionfinereference}).

